% v74 candidate: quantitative finite-group extension (all finite groups).
% This is a separately auditable strengthening of the preceding abelian appendix.
\section{Windows combined with arbitrary finite groups}
\label{app:window-group}

Here we remove commutativity from the finite-group hypothesis in
Appendix~\ref{app:window-abelian}.  The learning operator and the
window-only bound are unchanged.  All groups in this appendix are
fixed and finite.  The constants below are deliberately very large:
the result concerns polynomial dependence on the \emph{original}
grammar size and rule thickness, not a practical uniform bound in
the order of the group.

\begin{lemma}[A homogeneous block-order lemma]
\label{lem:group-block-order}
Fix positive integers $e,n$ with $n\ge3$.  For each
$t\in[e]$ and $i\in[n]$, let $(t,i)$ be a distinct point
on a line, with a total left-to-right order.  Suppose that
(i) the order of $(t,i)$ and $(u,i)$ is independent of $i$;
and (ii) for each ordered pair $(t,u)$, including $t=u$,
the order of $(t,i)$ and $(u,j)$ for $i<j$ is independent
of the choice of $i<j$.
Then the types at the points have a partition into classes,
each occupying a consecutive region of the line.  Every
class has one direction (increasing or decreasing in $i$);
inside its region the points occur in $n$ successive
\emph{chunks}, one per rank in that direction.  Each
chunk consists of exactly one point of each type in the
class (in a fixed within-chunk order).
\end{lemma}
\begin{proof}
By (ii) with $t=u$, each type's points are monotone.
Consider distinct types $t,u$.  Write
$X$ for $(t,i)<(u,j)$ when $i<j$, and $Y$ for
$(u,i)<(t,j)$ when $i<j$; both are uniform by (ii).
If both types increase, $X\wedge Y$ puts all points
of rank $i$ before all points of rank $j>i$.
The cases $X\wedge\neg Y$ and
$\neg X\wedge Y$ put all $t$-points before all
$u$-points, or conversely.  The remaining case
$\neg X\wedge\neg Y$ would give
$(t,1)<(t,2)<(u,1)<(u,2)<(t,1)$,
a contradiction.  If their directions are opposite,
$X\wedge Y$ and $\neg X\wedge\neg Y$
yield a cycle using indices $1,2,3$.
In the other two cases one type is entirely to
the left of the other: off-diagonal comparisons
follow from (ii), and comparisons at equal indices
follow from (i) and monotonicity.  Reverse indices
to cover the two-decreasing case.  Thus any two
types are either \emph{separated}, or have the
same direction and are \emph{interleaved by rank}.

Interleaving is transitive.  Indeed, if $t$ interleaves
with $u$ and $u$ with $v$, all three have the same
direction.  Reverse indices to assume increasing.
If $t$ and $v$ were separated with $t$ to the left,
we would have
$(v,1)<(u,2)<(t,3)<(v,1)$;
the other separation order gives the analogous
cycle.  Hence interleaving, with equality, partitions
the types into classes.  A type outside a class is
separated from each type inside it, and must lie
on the same side of all its types: otherwise two
interleaved types would themselves be separated.
The class regions are therefore consecutive.  The
rank comparisons give successive chunks.  Finally
the within-chunk order is independent of rank by (i).
\end{proof}

\begin{lemma}[Simultaneous finite-group interval shortening]
\label{lem:group-interval}
Let $A$ be any finite group.  Suppose $n$ indices each
contribute at most $f$ disjoint blocks to at most $f$
ordered output components, and each block has a prescribed
replacement block.  A selected index replaces \emph{all}
its blocks simultaneously.  The other material is fixed.
If the individual block destinations and within-index orders
are constant across the indices, and all cross-index block
order comparisons are constant for $i<j$, then
$n\ge\max\{3,|A|^f\}$ implies that a nonempty consecutive
interval of indices can be replaced without changing
any output component's value in $A$.
\end{lemma}
\begin{proof}
Apply Lemma~\ref{lem:group-block-order} separately to
the block types going to each output component.  Across
all output components there are $q\le f$ consecutive
class regions.  Consider such a region $K$, and enumerate
its chunks $1,\ldots,n$ from left to right.  Include all
unchanged material between the blocks in the region.
Let $P_p$ be the original group value of its prefix
ending after chunk $p$, and $Q_p$ the value of the
same prefix when chunks $1,\ldots,p$ have been
replaced; set $P_0=Q_0=1_A$.  Write
$\kappa^K_p=P_pQ_p^{-1}$.

Replacing just chunks $p+1,\ldots,p'$ gives the
following \emph{ordered} value for the entire region:
\[
 P_p\,Q_p^{-1}Q_{p'}\,P_{p'}^{-1}P_n .
\]
Indeed, $Q_p^{-1}Q_{p'}$ is exactly the replacement
value of the segment between the two prefix boundaries,
and $P_{p'}^{-1}P_n$ is its unchanged suffix.
Consequently, the region retains its original value
$P_n$ whenever $\kappa^K_p=\kappa^K_{p'}$.
This computation uses inverses, but never commutes
factors.

If class $K$ increases with the original index order,
its chunks for indices $(a,b]$ occupy $(a,b]$.
If it decreases, they occupy $(n-b,n-a]$.
For $j=0,\ldots,n$ define a state in $A^q$ by taking
$\kappa^K_j$ from each increasing class and
$\kappa^K_{n-j}$ from each decreasing class.
There are $n+1>|A|^q$ such states, so two,
at indices $a<b$, coincide.  Replacing precisely
indices $a+1,\ldots,b$ then preserves each
class region's value.  The rest of every output
component is unchanged, so all output group values
are preserved.
\end{proof}

\begin{lemma}[Uniform shortening of heavy antichains]
\label{lem:group-heavy-antichain}
Set
\[
 N_{A,f}=\max\{3,|A|^f\},\quad
 c_f=2^{f^2},\quad U_f=f^{f+1}f!,\quad
 M_{A,f}=U_f\,c_f^{\,c_f(N_{A,f}-1)+1}.
\]
Let $G$ be reduced and good, and let $g:\Sigma^*\to A$
be any homomorphism into a finite group.  In a terminal
derivation minimizing first total yield length and then
node count subject to its prescribed componentwise
$g$-value, every antichain of subtrees of total yield
length $>\bar\theta_G$ has fewer than $M_{A,f}$ nodes.
\end{lemma}
\begin{proof}
Every child component passes through the linear,
nondeleting templates to exactly one contiguous block
in a root component.  Blocks of incomparable subtrees
are disjoint.  Fix for each node of a heavy antichain
a replacement derivation of the \emph{same untyped}
nonterminal whose yield has length at most
$\bar\theta_G$; such a derivation exists by
Lemma~\ref{lem:thickness-local-bounds}(ii).
Its group values may differ from the original values.

An antichain node has at most $f$ blocks.  Its
\emph{unary pattern} records the number of blocks,
their root-component destinations, and their order
within any shared root component.  There are at
most $U_f$ such patterns.  First retain a set of
nodes having one common unary pattern.  For indices
$i<j$, color $\{i,j\}$ by all comparisons between
block types $(t,i)$ and $(u,j)$ situated in the
same root component, including $t=u$.  There are
at most $c_f$ colors.  The elementary multicolor
Ramsey bound
\[
 R_{c_f}(N)\le c_f^{\,c_f(N-1)+1}
\]
suffices: greedily select $c_f(N-1)+1$ successive
vertices, at each step retaining a largest
monochromatic set of later neighbors, and then
choose $N$ vertices whose selected outgoing color
is the same.  Each step preserves at least
$1/c_f$ of the available vertices (with integer
rounding).  Thus an antichain of at least
$M_{A,f}$ nodes contains $N=N_{A,f}$ indices
with constant unary pattern and constant
cross-index comparison color.

Within each root component their blocks meet
the hypotheses of
Lemma~\ref{lem:group-block-order}.
Lemma~\ref{lem:group-interval} therefore supplies
a nonempty interval of these selected indices
whose simultaneous replacement preserves the
\emph{whole root group vector}.  Every replaced
subtree was heavy but is replaced by one of length
at most $\bar\theta_G$, so the resulting valid
derivation has strictly smaller root yield length.
This contradicts its choice.
\end{proof}

\begin{theorem}[Short tuples with prescribed arbitrary finite-group values]
\label{thm:group-typed-shortest}
Let $G=(V,\Sigma,P,S,\fan)$ be a reduced good MCFG
of fan-out at most $f$ and maximum rule rank
$r_{\max}\ge1$.  Let $g:\Sigma^*\to A$ be a
homomorphism into \emph{any} finite group $A$.
For every $B\in V$ and every realizable
$\mathbf s\in A^{\fan(B)}$ there is
$\tuple u\in L_B(G)$ with
$g^{(\fan(B))}(\tuple u)=\mathbf s$ and
\[
 \|\tuple u\|_1\le
 \beta_g^{\mathrm{grp}}:=
 (r_{\max}+1)\bar\theta_G
 \max\{1,(M_{A,f}-1)|V||A|^f\}.
\]
In particular $\beta_g^{\mathrm{grp}}
=O_{A,f,r_{\max}}(\|G\|\bar\theta_G)$.
\end{theorem}
\begin{proof}
Choose a terminal derivation from $B$ having the
prescribed root vector, lexicographically minimizing
(total length, node count).  Label each node by
its nonterminal and componentwise group value.
An ancestor and a proper descendant cannot have
the same label: replacing the ancestor by the
descendant would preserve every ancestral group
vector (by the ordered evaluation of rule templates),
would not increase total length (by nondeletion),
and would decrease node count.  Hence each
root-to-leaf path has at most
$n_g=|V||A|^f$ nodes.

Call a node heavy when its subtree's total yield
length exceeds $\bar\theta_G$.  Heavy nodes are
ancestor-closed.  By
Lemma~\ref{lem:group-heavy-antichain},
their leaves number at most $M_{A,f}-1$;
all heavy nodes lie on paths to those leaves,
so their number is at most
$(M_{A,f}-1)n_g$.
If the root is light, its yield already has
length at most $\bar\theta_G$.
Otherwise, each heavy rule contributes at most
$\bar\theta_G$ explicit terminal symbols, and
each of its at most $r_{\max}$ light child
subtrees contributes at most $\bar\theta_G$.
Counting the heavy nodes and their maximal light
children gives the displayed upper bound.
\end{proof}

\begin{theorem}[Polynomial characteristic data for window--finite-group typing]
\label{thm:window-group-polydata}
Fix $f,r_{\max},k,\ell,\Sigma$, any finite group
$A$, a homomorphism $g:\Sigma^*\to A$, and a
fixed finite typing $h$ satisfying
$h\preceq(h_\Sigma^{k,\ell},g)$.
For every reduced good fan-out-$\le f$,
rank-$\le r_{\max}$ grammar $G$ of a target in
$\Cmcf{f}{h}$, the capped reconstruction operator
$\mathcal B^{\le r_{\max}}_{f,h}$ has a characteristic
sample such that
\[
 |\Wit(\widetilde G_0)|=O(\|G\|),\qquad
 \|\Wit(\widetilde G_0)\|_+
 =O_{f,r_{\max},k,\ell,h,g,\Sigma}
   (\|G\|^3\bar\theta_G).
\]
\end{theorem}
\begin{proof}
The proof of
Lemma~\ref{lem:window-abelian-short} uses
commutativity only through its input bound
Lemma~\ref{lem:group-typed-shortest}; its
skeleton replacement itself preserves the
group value of \emph{each omitted child component}
individually.  Replace that input by
Theorem~\ref{thm:group-typed-shortest}.
Specifically, retain the paths to the at most
$f(k+\ell+1)$ protected root terminal positions;
the skeleton has at most
$f(k+\ell+1)n_{\mathrm{typ}}$ nodes.
Each omitted child is replaced by a tuple of
the same underlying nonterminal with the same
componentwise $g$-vector and length at most
$\beta_g^{\mathrm{grp}}$.  Every replacement
component remains nonempty, so the marked
prefix--suffix positions remain unchanged and
the original long root components stay long.
Ordered evaluation of the unchanged child group
vectors preserves the root group vector
\emph{without commutativity}.  Thus the root
$h$-type is preserved and the new derivation
lifts to the trimmed full refinement.
The canonical tuples consequently have length
at most
\[
 f(k+\ell+1)n_{\mathrm{typ}}
 (\bar\theta_G+r_{\max}\beta_g^{\mathrm{grp}})
 =O(\|G\|^2\bar\theta_G),
\]
since $n_{\mathrm{typ}}\le|V||M|^f=O(\|G\|)$.
Lemma~\ref{lem:untyped-contexts} now bounds
each characteristic word by
$2\bar\theta_G+r_{\max}b_{\mathrm{can}}
=O(\|G\|^2\bar\theta_G)$, while
Proposition~\ref{prop:witness-cardinality}
bounds the number of words by $O(\|G\|)$.
The optional empty-word witness adds only a
constant to the encoded total.
\end{proof}

\begin{remark}
This strengthening does not assert an original-thickness
polynomial for arbitrary finite \emph{monoid} typings.
The explicit constant $M_{A,f}$ is a coarse
Ramsey bound.  The abelian argument of
Appendix~\ref{app:window-abelian} gives a substantially
smaller constant when $A$ is commutative.
\end{remark}
